\documentclass[11pt]{article}
\usepackage{hw-solutions-common}
\usepackage{amsmath,amssymb}

\begin{document}

\hwsolhead{Test 2}{Both tracks (PHYS 230 algebra / PHYS 235 calculus) share
Sections 1--3; Section 4 differs by track and is given separately below.
Constants used throughout: $\dfrac{1}{4\pi\varepsilon_0}=8.99\times10^{9}$
N$\cdot$m$^2$/C$^2$, $\varepsilon_0=8.85\times10^{-12}$ C$^2$/(N$\cdot$m$^2$),
$e=1.602\times10^{-19}$ C, $m_p=1.673\times10^{-27}$ kg.}

%============================================================
\problem{Section 1 --- P1: sign of $V$, direction of $\vec E$}{4}{$Q=+4.0$ nC
sits at the origin, $q=-1.0$ nC sits at $x=5.0$ cm. P1 is on the axis at
$x=4.0$ cm, 4.0 cm from $Q$ and 1.0 cm from $q$.}

\solve{Superpose the two point-charge potentials (dropping the common factor
$1/4\pi\varepsilon_0$ and units, as the hint allows):
\[
V_1 \propto \frac{Q}{r_Q}+\frac{q}{r_q} = \frac{4.0}{4.0}+\frac{-1.0}{1.0} = 1.0-1.0 =
\boxed{0}.
\]
\textbf{(C) zero.} For the field: $Q$ is positive and P1 is to its right, so
$\vec E_Q$ points away from $Q$, in $+\hat x$. $q$ is negative and P1 is to
its left, so $\vec E_q$ points \emph{toward} $q$ -- also $+\hat x$. The two
contributions add instead of fighting each other, so
\[
\boxed{\vec E_1 \text{ points in } +\hat x \text{ (due right, toward } q\text{)}}.
\]}

\review{P1 is the case where $V=0$ but $\vec E\neq 0$ -- the two potentials
cancel exactly (closer to the smaller-magnitude charge by exactly the ratio
that makes it work), but the two field contributions don't fight each other
at all, they add. $V=0$ and $\vec E=0$ are not the same condition, and this
point is designed to make that visible.}

%============================================================
\problem{Section 1 --- P2: sign of $V$, direction of $\vec E$}{4}{P2 is at
$(5.0,\,2.4)$ cm -- directly above $q$, 2.4 cm away from it, and
$\sqrt{5.0^2+2.4^2}=5.5$ cm from $Q$.}

\solve{\[
V_2 \propto \frac{4.0}{5.5}+\frac{-1.0}{2.4} = 0.73-0.42 = \boxed{+0.30}
\;\Rightarrow\; \textbf{(A) positive}.
\]
For the field, weight each unit vector by $1/r^2$. $\vec E_Q$ points away
from $Q$, along the line from $Q$ to P2 (up and to the right); $\vec E_q$
points straight down toward $q$ (P2 sits directly above it). Using
$ E\propto |charge|/r^2$:
\[
\vec E_{Q,2}\propto \frac{4.0}{5.5^2}(0.90,\,0.43) = (0.117,\,0.056),
\qquad
\vec E_{q,2}\propto \frac{1.0}{2.4^2}(0,\,-1) = (0,\,-0.174).
\]
Adding, $\vec E_2 \propto (0.117,\,-0.117)$ -- equal magnitude $+x$ and
$-y$ components, i.e.
\[
\boxed{\vec E_2 \text{ points } 45^\circ \text{ below } +\hat x \text{ (southeast)}}.
\]}

\review{The $2.4$ cm height isn't an arbitrary number -- it's exactly the
value that makes $Q$'s field and $q$'s field add up to a clean $45^\circ$
direction, so a numeric sanity check (rather than just a qualitative
``down-and-to-the-right'' guess) is the fastest way to confirm you have the
right compass direction.}

%============================================================
\problem{Section 1 --- P3: sign of $V$, direction of $\vec E$}{4}{P3 is on
the axis at $x=10$ cm, beyond $q$: 10 cm from $Q$, 5.0 cm from $q$.}

\solve{\[
V_3 \propto \frac{4.0}{10}+\frac{-1.0}{5.0} = 0.40-0.20 = \boxed{+0.20}
\;\Rightarrow\; \textbf{(A) positive}.
\]
For the field: $\vec E_Q \propto Q/r_Q^2 = 4.0/10^2 = 0.040$, pointing away
from $Q$ (i.e.\ $+\hat x$, since P3 is to the right of $Q$). $\vec E_q
\propto |q|/r_q^2 = 1.0/5.0^2 = 0.040$, pointing \emph{toward} $q$ (i.e.\
$-\hat x$, since P3 is to the right of $q$). Equal magnitude, opposite
direction:
\[
\boxed{\vec E_3 = 0}.
\]}

\review{P3 is the mirror image of P1: there, $V$ canceled but $\vec E$
didn't; here $\vec E$ cancels exactly (by design: $4.0/10^2$ and
$1.0/5.0^2$ both equal $0.040$) while $V$ stays comfortably positive. Together
the three points cover all three cases a student needs to tell apart:
$V=0,\vec E\neq0$ (P1); both nonzero (P2); $V\neq0,\vec E=0$ (P3).}

%============================================================
\problem{Section 2 --- Part 1: the dipole moment}{5}{Charges of magnitude
5 nC, separated by $6\ \mu$m, with $\vec E = 20$ N/C $\hat y$. The dipole
moment points $60^\circ$ below the $x$-axis (``Orientation A'').}

\solve{The dipole moment has magnitude
\[
p = qd = (5.0\times10^{-9}\text{ C})(6.0\times10^{-6}\text{ m}) =
\boxed{3.0\times10^{-14}\text{ C}\cdot\text{m}},
\]
pointing (by definition, from $-q$ to $+q$) $60^\circ$ below the $+x$-axis,
i.e.
\[
\vec p = p(\cos(-60^\circ),\,\sin(-60^\circ),\,0) = (1.5,\,-2.6,\,0)\times10^{-14}\text{
C}\cdot\text{m}.
\]
\textbf{Drawing:} the $+x$-axis horizontal, $+y$ vertical. $\vec E$ is a
vertical arrow pointing straight up ($+\hat y$). $\vec p$ is a separate arrow
starting at the dipole's center, $60^\circ$ \emph{below} the $+x$-axis (down
and to the right, into the fourth quadrant) -- \emph{not} aligned with
$\vec E$ and not symmetric about either axis. The negative charge sits at
the tail of $\vec p$, the positive charge at the head, $6\ \mu$m apart along
that same line.}

\review{A common error here is drawing $p$ at $60^\circ$ \emph{above} the
axis, or measuring the angle from $\vec E$ instead of from the $x$-axis --
read the setup's wording carefully (as the problem itself warns): the angle
given is always relative to the axis stated, not to the field.}

%============================================================
\problem{Section 2 --- Part 2: torque}{4}{Same dipole and field as above.}

\solve{$\vec\tau = \vec p\times\vec E$. Since $\vec p=(p_x,p_y,0)$ and
$\vec E=(0,E_y,0)$ (purely $\hat y$),
\[
\vec\tau = (0,\,0,\,p_xE_y) = (0,\,0,\,(1.5\times10^{-14})(20)) =
\boxed{(0,\,0,\,3.0\times10^{-13})\text{ N}\cdot\text{m}},
\]
i.e.\ $3.0\times10^{-13}$ N$\cdot$m in $+\hat z$ (out of the page).}

\review{The torque vector's direction is the rotation \emph{axis}, found by
the right-hand rule, not a push in some in-plane direction. $+\hat z$ means:
curl the right hand's fingers from $\vec p$ toward $\vec E$ through the
\emph{smaller} angle between them, and the thumb points out of the page --
confirming the dipole is being spun counterclockwise (as drawn, with $+x$
right and $+y$ up), rotating $\vec p$ up and to the left, toward alignment
with $\vec E$.}

%============================================================
\problem{Section 2 --- Part 3: potential energy, Orientation A vs.\ B}{6}{Same
dipole and field. Orientation B is whichever orientation has the lowest PE.}

\solve{$U=-\vec p\cdot\vec E=-pE\cos\theta$, where $\theta$ is the angle
between $\vec p$ and $\vec E$. This is minimized when $\cos\theta=1$, i.e.\
$\vec p$ points the \emph{same direction as} $\vec E$ -- straight up,
$+\hat y$. \textbf{Orientation B: $\vec p$ parallel to $\vec E$ (pointing
straight up).}

In Orientation A, $\vec p$ is at $-60^\circ$ and $\vec E$ is at $+90^\circ$
(both measured from $+x$), so $\theta_A = 90^\circ-(-60^\circ)=150^\circ$:
\[
U_A = -pE\cos150^\circ = -(3.0\times10^{-14})(20)(-0.866) = +5.2\times10^{-13}\text{ J}.
\]
In Orientation B, $\theta_B=0$:
\[
U_B = -pE\cos0^\circ = -(3.0\times10^{-14})(20) = -6.0\times10^{-13}\text{ J}.
\]
So Orientation A has
\[
U_A-U_B = 5.2\times10^{-13}-(-6.0\times10^{-13}) = \boxed{1.1\times10^{-12}\text{
J \emph{more} than Orientation B}}.
\]}

\review{$U_A-U_B = pE(1-\cos\theta_A)$ ranges from $0$ (if A already equals
B) to $2pE=1.2\times10^{-12}$ J (if A is the \emph{worst} orientation,
anti-aligned with $\vec E$). $150^\circ$ is close to that worst case
($180^\circ$), so it makes sense that $U_A-U_B$ comes out close to, but
just under, that $2pE$ ceiling.}

%============================================================
\problem{Section 3 --- parts 1 \& 2: two charged rods}{16}{\emph{PENDING --
needs the figure.} The source file references an image, \texttt{rods.png},
showing where points A, B, and C sit relative to the two rods (B is given as
halfway between them, but A and C's exact positions/distances aren't
recoverable from the \LaTeX{} source alone). Once I have that figure (or
you just tell me how far A and C sit from their nearest rod, in terms of
$d$), I'll fill in both parts -- same verified-numerically approach as the
rest of this set. The method itself is straightforward superposition of
$E=\lambda/2\pi\varepsilon_0 r$ from each rod (with sign/direction from each
rod's charge), same pattern as Written HW 3's two-sheets problem.}

%============================================================
\begin{center}{\LARGE\bfseries Section 4 --- PHYS 235 (Calculus track): Gauss's Law}\end{center}
\vspace{2mm}

\problem{Part 1: enclosed charge}{8}{A cube with 3.0 cm sides. The field is
purely in $\hat y$ everywhere. At the top face, $\vec E=\langle
0,-34,0\rangle$ N/C; at the bottom face, $\vec E=\langle 0,+20,0\rangle$
N/C.}

\solve{Since $\vec E$ has no $x$ or $z$ component, the flux through all four
side faces is zero (their outward normals are $\pm\hat x$ or $\pm\hat z$,
perpendicular to $\vec E$). Only the top and bottom faces contribute. With
$A=(0.030\text{ m})^2 = 9.0\times10^{-4}$ m$^2$:

\emph{Top face}, outward normal $+\hat y$:
$\Phi_{\text{top}} = \vec E_{\text{top}}\cdot\hat n\,A = (-34)(+1)A = -34A$.

\emph{Bottom face}, outward normal $-\hat y$:
$\Phi_{\text{bottom}} = \vec E_{\text{bottom}}\cdot\hat n\,A = (+20)(-1)A = -20A$.

\[
\Phi_{\text{total}} = -34A-20A = -54A = -54(9.0\times10^{-4}) = -0.0486\text{
N}\cdot\text{m}^2/\text{C}.
\]
By Gauss's law,
\[
Q_{\text{enc}} = \varepsilon_0\Phi_{\text{total}} =
(8.85\times10^{-12})(-0.0486) = \boxed{-4.3\times10^{-13}\text{ C} \approx
-0.43\text{ pC}}.
\]
A drawing should show the cube with $\vec E$ arrows pointing \emph{down}
into the top face (since $\vec E_{\text{top}}$ points in $-\hat y$, into the
cube from above) and pointing \emph{up} into the bottom face (since
$\vec E_{\text{bottom}}$ points in $+\hat y$, into the cube from below) --
field lines converging into the cube from both the top and the bottom,
consistent with a net negative charge trapped inside.}

\review{The two faces don't need to have equal field strength for this to
work -- Gauss's law only cares about the net flux. Here $34$ and $20$ N/C
aren't equal, so the two faces don't exactly cancel, and that imbalance
(both pointing inward, but by different amounts) is exactly what leaves a
nonzero, negative $Q_{\text{enc}}$.}

\problem{Part 2: what Gauss's law can and can't tell you}{4}{Same cube.}

\solve{Gauss's law only ever gives the \emph{total} charge enclosed by a
surface -- it has no way to say \emph{where} inside the cube that charge
actually sits. The $-0.43$ pC found above could be a thin sheet right at the
cube's center, spread uniformly through the whole volume, concentrated near
one face, or any other distribution; this calculation can't distinguish
between them, and nothing in the given field values (just two numbers, at
the top and bottom) provides enough information to localize it further.

For charge \emph{outside} the cube: this particular Gaussian surface (the
cube itself) tells us nothing about it. To learn anything about charge
outside, you'd need to draw a \emph{different} Gaussian surface that
encloses some of that outside region, together with field information on
that new surface -- neither of which this problem gives us. The one thing
we \emph{can} say is that whatever is outside is consistent with the field
being purely $\hat y$-directed everywhere (not just inside the cube): that's
a strong symmetry statement (reminiscent of infinite charged sheets, where
the field depends only on height, never on the horizontal position), but it
still doesn't pin down a location or an amount.}

\review{This pair of questions is really testing the same idea from two
sides: Gauss's law is a statement about a \emph{surface integral} of the
field, which only ever recovers the enclosed charge's \emph{total}, not its
spatial distribution, and only for the one surface you chose to apply it to.}

%============================================================
\begin{center}{\LARGE\bfseries Section 4 --- PHYS 230 (Algebra track): Using Electric Potential}\end{center}
\vspace{2mm}

\problem{Part 1: does the proton reach point 2?}{6}{A proton passes point 1
at $v_1=50{,}000$ m/s, crossing from the $30$ V equipotential to the $-10$ V
equipotential at point 2.}

\solve{Energy conservation for a charge $q$ moving through a potential
difference: $\frac12mv_1^2+qV_1 = \frac12mv_2^2+qV_2$, so
\[
\frac12mv_2^2 = \frac12mv_1^2 + q(V_1-V_2).
\]
For the proton, $q=+e$, and $V_1-V_2 = 30-(-10)=40$ V, so the proton is
moving from high to low potential -- a positive charge ``sliding downhill,''
which \emph{speeds it up}:
\[
\tfrac12mv_1^2 = \tfrac12(1.673\times10^{-27})(50{,}000)^2 = 2.09\times10^{-18}\text{ J},
\]
\[
q(V_1-V_2) = (1.602\times10^{-19})(40) = 6.41\times10^{-18}\text{ J},
\]
\[
\tfrac12mv_2^2 = 2.09\times10^{-18}+6.41\times10^{-18} = 8.50\times10^{-18}\text{ J},
\]
\[
v_2 = \sqrt{\frac{2(8.50\times10^{-18})}{1.673\times10^{-27}}} =
\boxed{1.0\times10^{5}\text{ m/s. Yes, it reaches point 2.}}
\]}

\review{A proton moving from higher to lower potential is exactly like a
ball rolling downhill -- it's expected to speed up, and it does (from
$5.0\times10^4$ to $1.0\times10^5$ m/s, roughly double).}

\problem{Part 2: does the antiproton reach point 2?}{6}{An antiproton (same
mass, charge $-e$) passes point 1 with the \emph{same} speed, $50{,}000$
m/s, along the same path.}

\solve{Same equation, but now $q=-e$:
\[
\tfrac12mv_2^2 = \tfrac12mv_1^2 + q(V_1-V_2) = 2.09\times10^{-18} +
(-1.602\times10^{-19})(40),
\]
\[
\tfrac12mv_2^2 = 2.09\times10^{-18}-6.41\times10^{-18} =
-4.32\times10^{-18}\text{ J}.
\]
That's negative, which is impossible for $\tfrac12mv_2^2$ -- \textbf{the
antiproton does not reach point 2.} It's a negative charge moving from high
to low potential, which is ``uphill'' for it, so it slows down and turns
around before getting there. To find where: set the kinetic energy to zero
at the turning point and solve for $V_{\text{turn}}$,
\[
\tfrac12mv_1^2 + qV_1 = 0 + qV_{\text{turn}}
\;\Rightarrow\;
V_{\text{turn}} = V_1 + \frac{\tfrac12mv_1^2}{q} = 30 +
\frac{2.09\times10^{-18}}{-1.602\times10^{-19}} = 30-13.1 =
\boxed{+17\text{ V}}.
\]}

\review{$17$ V sits between the $30$ V starting point and the $-10$ V target
-- so the antiproton makes it about halfway there before running out of
kinetic energy and turning back. Same mass, same starting speed, opposite
charge, opposite outcome: this pair of problems is designed to make that
contrast (speeds up vs.\ turns around) as stark as possible.}

\end{document}
